%!TeX root=csp.tex
\section{What is being proved}\label{sec:target}

\subsection{The theorem}

Fix a finite set $A$ and a finite set $\Gamma$ of relations on $A$, the
\emph{constraint language}. An instance of $\CSP(\Gamma)$ is a conjunction
$R_{1}(\bar x_{1})\wedge\dots\wedge R_{s}(\bar x_{s})$ with each
$R_{i}\in\Gamma$; the question is whether it is satisfiable.

\begin{theorem}[Dichotomy; Bulatov, Zhuk]\label{thm:dichotomy}
If some weak near-unanimity operation preserves $\Gamma$, then $\CSP(\Gamma)$ is
solvable in polynomial time. Otherwise $\CSP(\Gamma)$ is \textup{NP}-complete.
\end{theorem}

Both halves are conditional on a model of computation, and this is not a
formality that a formalization may wave through. ``Solvable in polynomial time''
quantifies over an encoding of instances as strings, a machine, and a
polynomial bound on its step count; ``\textup{NP}-complete'' additionally
quantifies over all problems in \textup{NP} and over polynomial-time many-one
reductions. Mathlib supplies Turing machines
(\texttt{Mathlib/Computability/TuringMachine/}) and a notion of computability,
but it has \emph{no} complexity class, no polynomial-time predicate, and no
notion of polynomial-time reduction.

\subsection{Where the work actually is}\label{sec:program-axis}

An earlier draft of this section concluded from that gap that the
complexity-theoretic half should be deferred because building it would dominate
the project. That inference was wrong, and it is worth recording why, because
the corrected version is what organises the rest of the document.

The complexity vocabulary itself --- languages, a cost model, $\mathrm P$,
$\mathrm{NP}$ by verifiers, $\le_{p}$, \textup{NP}-hardness,
\textup{NP}-completeness, and a generic \textup{NP}-hard problem --- has since
been built, in about $850$ lines. That is consistent with the two existing
Cook--Levin formalizations, whose class layers are $702$ lines (Isabelle) and
$992$ lines (Coq) against totals of $56\,514$ and $\approx 16\,500$. Defining
the classes is cheap in every development that has done it.

What is expensive is not complexity theory. It is any statement that requires
one to \emph{exhibit a program}. In a formalization, ``$f$ is computable in
polynomial time'' is not a property one observes of a mathematical function; it
is an existential over programs, discharged only by writing one and proving a
bound on its step count. The dividing line runs along that axis, and it cuts
across the algebra/complexity split rather than agreeing with it:

\begin{center}
\begin{tabular}{@{}lll@{}}
\hline
Target & Exhibit a program? & Cost \\
\hline
\textup{(T0)} the algebraic core & no & the mathematics \\
\textup{(T1)} $\textsc{Solve}$ is correct & no (but see Hazard~\ref{haz:t1-computable}) & moderate \\
\textup{(T2)} $\textsc{Solve}$ is polynomial & \textbf{yes --- the largest one here} & \textbf{large} \\
\textup{(H0)} $\Gamma$ pp-interprets \textsc{Nae} & no & the mathematics \\
\textup{(H1)} pp-interpretation $\Rightarrow\le_{p}$ & yes, a syntactic substitution & small \\
\hline
\end{tabular}
\end{center}

So the document still separates the machine-free statements from the rest, but
for a sharper reason than before, and with a different verdict on \textup{(T2)}.

\subsection{The layered targets}

Write $\mathbf A=(A;w)$ for a finite idempotent algebra whose basic operation
$w$ is a special weak near-unanimity operation, and $\Vn$ for the class of such
algebras with $w$ of arity $n$ (Convention~\ref{conv:standing}).

\paragraph{Tractable half.} Three layers.

\begin{description}\tightlist
\item[\textup{(T0)} \emph{The algebraic core.}] The three theorems of
  \S\ref{sec:main-statements} --- the existence of a strong subuniverse or a
  linear congruence, the safety of a reduction to a strong subuniverse, and the
  codimension-one theorem --- together with everything they rest on. These are
  statements about finite algebras, congruences, and CSP instances regarded as
  finite combinatorial objects. No computation appears in them.
\item[\textup{(T1)} \emph{Correctness of the algorithm.}] Zhuk's algorithm
  written as a total function $\textsc{Solve}$ from instances to Booleans, and
  the proposition
  \[
    \textsc{Solve}(\mathcal I)=\texttt{true}\iff
    \mathcal I\ \text{has a solution}.
  \]
  This is a statement about a recursive function, and it is provable from
  \textup{(T0)} with no cost model. It is the honest formal content of ``the
  algorithm works'' --- \emph{provided} $\textsc{Solve}$ actually computes;
  see Hazard~\ref{haz:t1-computable}.
\item[\textup{(T2)} \emph{The complexity wrapper.}] That $\textsc{Solve}$ runs
  in polynomial time in the size of the instance, for fixed $\Gamma$. This is
  the one place a cost model is needed on this side, and it is a large
  undertaking rather than a wrapper: see Hazard~\ref{haz:t2-is-not-small}.
\end{description}

\paragraph{Hard half.} Two layers.

\begin{description}\tightlist
\item[\textup{(H0)} \emph{The algebraic core.}] If no weak near-unanimity
  operation preserves $\Gamma$ then $\Gamma$ primitively positively interprets
  every finite constraint language --- in particular one whose CSP is
  \textup{NP}-hard. This is again pure algebra: the Galois connection between
  polymorphisms and primitive positive definability, the Maróti--McKenzie
  theorem that a Taylor term yields a weak near-unanimity term, and the
  Bulatov--Jeavons--Krokhin reduction.
\item[\textup{(H1)} \emph{The complexity wrapper.}] That a primitive positive
  interpretation yields a polynomial-time many-one reduction, and hence
  \textup{NP}-hardness. This does need a program, but the program performs a
  syntactic substitution of formulas whose size is linear in the instance, so
  it is genuinely small.
\end{description}

The claim this document makes is that \textup{(T0)}, \textup{(T1)} and
\textup{(H0)} are the mathematics and are what the literature actually proves.
\textup{(H1)} is a small addition on top. \textup{(T2)} is neither small nor
mathematics, but it is not \emph{unrelated} to the dichotomy either --- it is
the assertion that the algorithm one has just proved correct is also fast, and
it is the second-largest component of the project. The blueprint is written for
\textup{(T0)}, \textup{(T1)} and \textup{(H0)};
\S\ref{sec:complexity-layer} states \textup{(T2)} and \textup{(H1)} precisely
and says what building them costs.

\begin{hazard}[\textup{(T1)} is vacuous unless $\textsc{Solve}$ computes]
\label{haz:t1-computable}
``$\textsc{Solve}$ written as a total function'' hides a requirement that no
cost model would catch. If $\textsc{Solve}$ is defined using classical choice
--- for instance by deciding ``does $D_{x}$ have a nonempty proper binary
absorbing subuniverse?'' with \texttt{Classical.dec} --- then
$\textsc{Solve}(\mathcal I)=\texttt{true}\iff\mathcal I$ has a solution is a
true theorem that says nothing algorithmic, because $\textsc{Solve}$ does not
reduce. To mean what it appears to mean, \textup{(T1)} needs every predicate
the algorithm branches on to carry a genuine decision procedure.

That is a real obligation with a locatable cost. Absorption is defined by an
existential over \emph{terms} --- an infinite type --- so
``$C\subt{\TBA}B$'' is not decidable on its face. It becomes decidable through
the bound that a witness of arity at most $|A|^{|A|}$ suffices, which is a
theorem and not a definition. Hazard~\ref{haz:algorithm-formalization}(b)
records that the bound exists; what it does not record, and what matters here,
is that \textup{(T1)} is empty without it. The same applies to the searches
over congruences and over subsets.
\end{hazard}

\begin{hazard}[\textup{(T2)} is not a wrapper]\label{haz:t2-is-not-small}
Two things make \textup{(T2)} much larger than the phrase ``complexity
wrapper'' suggests.

\emph{The program is enormous.} \textup{(T2)} requires $\textsc{Solve}$ ---
\textsc{ForceConsistency}, the search for a strong subuniverse,
\textsc{SolveLinear}, and the mutual recursion between them --- to be expressed
in a cost model, with a proved polynomial bound. Every existing measurement
says this kind of work dominates: it is the $50\,000$ lines of Turing-machine
engineering in the Isabelle Cook--Levin development, against $702$ for the
classes themselves.

\emph{The statement is non-uniform.} Because $\Gamma$ is fixed, the polynomial
may depend on $\Gamma$, and the searches over subsets of $A$, over congruences
on $A$, and over term operations of bounded arity are all constant-time. But
``constant'' here means the search is compiled away into a program whose
\emph{size} depends on $|A|$. So $\textsc{Solve}$ is not one program: it is a
family $\textsc{Solve}_{\Gamma}$ with a polynomial $p_{\Gamma}$ for each
$\Gamma$, and \textup{(T2)} must be stated in that form. Attempting a bound
uniform in $\Gamma$ would be proving something else, and something false.
\end{hazard}

\begin{remark}
The separation is not a dodge. Layer \textup{(T1)} is strictly stronger than
``\,$\CSP(\Gamma)$ is decidable'', which is trivial for a finite template: it
says that one particular, highly non-obvious algorithm --- the one whose
existence is the content of the dichotomy --- is correct. Everything that makes
the dichotomy hard lives in \textup{(T0)}. Conversely, no part of
\textup{(T0)}--\textup{(T1)} becomes easier if a cost model is available.
\end{remark}
